%%%%%%%%%%%%%%%%%%%%%%%%%%%%%%%%%%%%%%%%%%%%%%%%%%%%%%%%
\subsection{Mass Model for the Full CBC Population}\label{appendix:PDB}
%%%%%%%%%%%%%%%%%%%%%%%%%%%%%%%%%%%%%%%%%%%%%%%%%%%%%%%%

Our \parametricAdj \PDB mass model is based on \textsc{Power Law+Dip}~\citep{Fishbach:2020ryj}, \textsc{Broken Power Law + Dip}~\citep{Farah:2021qom}, and \textsc{MultiPDB}~\citep{Mali:2024wpq}.
It is also similar to the \textsc{Power Law+Dip+Break} model~\citep{KAGRA:2021duu}, but in our \PDB model we now include additional structure.
We allow for the possibility of an upper mass gap, as well as separate pairing functions for \ac{NS}-containing binaries and \acp{BBH}.
Additionally, we include explicit peaks at low and mid-range \ac{BH} masses to capture the $\sim 9$--$10\,\Msun$ and $\sim 30$--$40\, \Msun$ features found by the \nonparametrics.
As discussed in Section~\ref{sec:all_masses}, we find that the addition of a low-mass \ac{BH} peak eliminates the need to explicitly model a gap between \acp{NS} and \acp{BH} masses.

Our mass model is parameterized as
\begin{equation}
    \pi(m_1, m_2|\PEhyperparam) \propto \pi_m(m_1|\PEhyperparam)\pi_m(m_2|\PEhyperparam)f(m_1, m_2)\Theta(m_2 < m_1)
\label{eq: pdb joint mass model}
\end{equation}
where the one-dimensional mass distribution $\pi_m(m|\PEhyperparam)$ is given by
\begin{align}
\pi_m&(m|\PEhyperparam) =  \left[1 + c_1 \mathcal{N}_{[m_{\text{min}}, m_{\text{max}}]}(m|\mu_1, \sigma_1) + c_2 \mathcal{N}_{[m_{\text{min}}, m_{\text{max}}]}(m|\mu_2, \sigma_2)\right] \nonumber \\
& \times n_1(m|m_\text{max}^\text{NS}, m_\text{min}^\text{BH}, \eta_\text{max}^\text{NS}, \eta_\text{min}^\text{BH}, A) \, n_2(m|m_\text{min}^\text{UMG}, m_\text{max}^\text{UMG}, \eta_\text{min}^\text{UMG}, \eta_\text{max}^\text{UMG}, A_2) \nonumber \\
& \times h(m|m_\text{min}^\text{NS}, \eta_\text{min}^\text{NS}) \,l(m|m_\text{max}^\text{BH}, \eta_\text{max}^\text{BH}) 
\begin{cases}
m^{\alpha_1} & m < m_\text{max}^\text{NS} \\
m^{\alpha_{\text{dip}}} [m_\text{max}^\text{NS}]^{\alpha_1 - \alpha_{\text{dip}}} &  m \in [m_\text{max}^\text{NS}, m_\text{min}^\text{BH}) \\
m^{\alpha_2} [m_\text{max}^\text{NS}]^{\alpha_1 - \alpha_{\text{dip}}} [m_\text{min}^\text{BH}]^{\alpha_{\text{dip}} - \alpha_2} & m \geq m_\text{min}^\text{BH}.
\end{cases}
\label{eq: pdb mass model}
\end{align}

This $\pi_m(m|\PEhyperparam)$ represents a universal mass function to describe the primary and secondary mass distributions. Note that the marginal mass distribution is different from the universal mass distribution due to the pairing formalism.
$\mathcal{N}_{[a,b]}(\mu, \sigma)$ represents a truncated normal distribution over $[a,b]$ with location and width parameters $\mu$ and $\sigma$. 
The high-pass, low-pass and notch functions are defined as follows:

\begin{align}
    \,l(m|m_\text{max}^\text{BH}, \eta_\text{max}^\text{BH}) &= \left[1+\left(\frac{m}{m_\text{max}^\text{BH}}\right)^{\eta_\text{max}^\text{BH}}\right]^{-1},\\
    h(m|m_\text{min}^\text{NS}, \eta_\text{min}^\text{NS}) &= 1 - l(m|m_\text{min}^\text{NS}, \eta_\text{min}^\text{NS}),\\
    n_1(m|m_\text{max}^\text{NS}, m_\text{min}^\text{BH}, \eta_\text{max}^\text{NS}, \eta_\text{min}^\text{BH}, A) &= 1 - A \,l(m|m_\text{max}^\text{NS}, \eta_\text{max}^\text{NS}) \,h(m|m_\text{min}^\text{BH}, \eta_\text{min}^\text{BH}),\\
	n_2(m|m_\text{min}^\text{UMG}, m_\text{max}^\text{UMG}, \eta_\text{min}^\text{UMG}, \eta_\text{max}^\text{UMG}, A_2) &= 1 - A_2 \,l(m|m_\text{min}^\text{UMG}, \eta_\text{min}^\text{UMG}) \,h(m|m_\text{max}^\text{UMG}, \eta_\text{max}^\text{UMG}).\\
\end{align}

The pairing function
\begin{equation}
f(m_1,m_2|\beta_{\rm BH}, \beta_{\rm NS}) = \begin{cases}
    \left(\dfrac{m_2}{m_1}\right)^{\beta_{1}} & {\rm if} \; m_2 < 5\,\Msun \\
    \left(\dfrac{m_2}{m_1}\right)^{\beta_{2}} & {\rm if} \; m_2 > 5\,\Msun \\
\end{cases}
\end{equation}
controls how much merging binaries favor/disfavor equal masses. We allow for alternative pairing for binaries with very light secondary masses (\acp{NSBH} or \ac{BBH} with the secondary component in the lower-mass gap, e.g., GW190814 \citealt{LIGOScientific:2020zkf}).
We show the priors and describe the parameters of the \PDB model in Table~\ref{tab:parameters_pdb}.

\begin{deluxetable}{ccccc}
\tablecaption{Summary of \PDB model parameters and priors.}
\tablehead{
\colhead{Category} & \colhead{Parameter} & \colhead{Unit} & \colhead{Description} & \colhead{Prior}
}
\label{tab:parameters_pdb}
\startdata
Pairing Function & $\beta_1$ & -- & Spectral index below $5 \, \Msun$ & U($-2, 3$) \\
& $\beta_2$ & -- & Spectral index above $5\, \Msun$ & U($-2, 7$) \\
\colrule
Broken Power-Law & $\alpha_1$ & -- & Powerlaw below $m_{\text{NS max}}$ & U($-10, 2$) \\
& $\alpha_{\text{dip}}$ & -- & Powerlaw between  & U($-3, 2$) \\
&  &  & $m_{\text{NS max}}$ and $m_{\text{BH min}}$ &  \\
& $\alpha_2$ & -- & Powerlaw above $m_{\text{BH min}}$ & U($-3, 2$) \\
& $m_{\text{brk}}$ & $\Msun$ & Break point between $\alpha_1$ and $\alpha_2$ & $5$ \\
\colrule
Highpass Filter & $m_{\text{NS min}}$ & $\Msun$ & Low-mass roll-off & U($1, 1.4$) \\
& $\eta_{\text{min}}$ & -- & Sharpness at $m_{\text{NS min}}$ & $50$ \\
\colrule
Lowpass Filter & $m_\text{max}^\text{BH}$ & $\Msun$ & High-mass roll-off & U($60, 200$) \\
& $\eta_{\text{max}}$ & -- & Sharpness at $m_\text{max}^\text{BH}$ & U($-4, 12$) \\
\colrule
Low-Mass Notch & $m_{\text{NS max}}$ & $\Msun$ & Lower notch edge & U($1.4, 5$) \\
& $\eta^{\text{low}}_1$ & -- & Sharpness at $m_{\text{NS max}}$ & $50$ \\
& $m_{\text{BH min}}$ & $\Msun$ & Upper notch edge & U($5, 9$) \\
& $\eta^{\text{high}}_1$ & -- & Sharpness at $m_{\text{BH min}}$ & $50$ \\
& $A_1$ & -- & Notch depth & $0$ \\
\colrule
High-Mass Notch & $m_\text{min}^\text{UMG}$ & $\Msun$ & Lower notch edge & U($30, 90$) \\
& $\eta^{\text{low}}_2$ & -- & Sharpness at $m_\text{min}^\text{UMG}$ & $30$ \\
& $m_\text{max}^\text{UMG}$ & $\Msun$ & Upper notch edge & U($60, 150$) \\
& $\eta^{\text{high}}_2$ & -- & Sharpness at $m_\text{max}^\text{UMG}$ & $30$ \\
& $A_2$ & -- & Depth of high-mass notch & U($0, 1$) \\
\colrule
Low-Mass Peak & $\mu^{\text{peak}}_2$ & $\Msun$ & Peak location & U($6, 12$) \\
& $\sigma^{\text{peak}}_2$ & $\Msun$ & Peak width & U($0, 5$) \\
& $c_2$ & -- & Peak height & U($0, 500$) \\
\colrule
High-Mass Peak & $\mu^{\text{peak}}_1$ & $\Msun$ & Peak location & U($17, 50$) \\
& $\sigma^{\text{peak}}_1$ & $\Msun$ & Peak width & U($4, 20$) \\
& $c_1$ & -- & Peak height & U($0, 1000$) \\
\enddata
\tablecomments{U($x$, $y$) denotes a Uniform prior between $x$ and $y$.}
\end{deluxetable}

%%%%%%%%%%%%%%%%%%%%%%%%%%%%%%%%%%%%%%%%%%%%%%%%%%%%%%%%
\subsection{Neutron Star Mass Models}\label{appendix:NSmodels}
%%%%%%%%%%%%%%%%%%%%%%%%%%%%%%%%%%%%%%%%%%%%%%%%%%%%%%%%

\begin{deluxetable}{cccc}
\tablecaption{Summary of \textsc{Power} and \textsc{Peak} \ac{NS} mass model parameters.}
\tablehead{
\colhead{Parameter} & \colhead{Unit} & \colhead{Description} & \colhead{Prior}
}
\label{tab:parameters_ns}
\startdata
$\alpha$ & -- & Spectral index for the power-law in the  & U($-15$, $5$) \\
&  & \textsc{Power} NS mass distribution. &  \\
$m_{\rm min}$ & $\Msun$ & Minimum mass of the \ac{NS} mass distribution. & U($1.0$, $1.5$) \\
$m_{\rm max}$ & $\Msun$ & Maximum mass of the \ac{NS} mass distribution. & U($1.5$, $3.0$) \\
$\mu$ & $\Msun$ & Location of the Gaussian peak in the  & U($1.0$, $3.0$) \\
 & & \textsc{Peak} \ac{NS} mass distribution. &  \\
$\sigma$ & $\Msun$ & Width of the Gaussian peak in the & U($0.01$, $2.00$) \\
 &  & \textsc{Peak} \ac{NS} mass distribution. & \\
\enddata
\tablecomments{U($x$, $y$) denotes a Uniform prior between $x$ and $y$.}
\end{deluxetable}

\noindent Following previous work~\citep{Landry:2021hvl,KAGRA:2021duu},
the mass distribution of \ac{NS}-containing events is modeled as
\begin{equation}
\pi(m_1, m_2 | \PEhyperparam) \propto
\begin{cases}
\pi(m_1| \PEhyperparam) \, \pi(m_2| \PEhyperparam) & \text{if BNS,} \\
\mathrm{U}(3\,\Msun, 60\,\Msun) \, \pi(m_2 | \PEhyperparam) & \text{if NSBH},
\end{cases}
\end{equation}
To model $\pi(m|\PEhyperparam)$, we use either of the following models
\begin{enumerate}
	\item \textsc{Power} model:
\begin{equation}
	\pi(m|\PEhyperparam) \propto
	\begin{cases}
		m^{\alpha} & \text{if } m_{\min} \leq m \leq m_{\max}, \\
		0 & \text{otherwise.}
	\end{cases}
\end{equation}

\item \textsc{Peak} model:
\begin{equation}
	\pi(m|\PEhyperparam) \propto
	\begin{cases}
		\exp\left[ -\dfrac{(m - \mu)^2}{2\sigma^2} \right] & \text{if } m_{\min} \leq m \leq m_{\max}, \\
		0 & \text{otherwise.}
	\end{cases}
\end{equation}

See Table~\ref{tab:parameters_ns} for a description of the parameters used in the model and the corresponding prior ranges.

\end{enumerate}


%%%%%%%%%%%%%%%%%%%%%%%%%%%%%%%%%%%%%%%%%%%%%%%%%%%%%%%%
\subsection{Binary Black Hole Mass Models}\label{appendix:BrokenPLTwoPeaks}
%%%%%%%%%%%%%%%%%%%%%%%%%%%%%%%%%%%%%%%%%%%%%%%%%%%%%%%%

\BrokenPLTwoPeaks : The fiducial \ac{BBH} mass model is a mixture between a broken power law and two left-truncated Gaussian peaks, with low mass tapering applied to the full distribution. The broken power law is given by

\begin{eqnarray}
	p_\text{BP}(m_1 | \alpha_1, \alpha_2, m_\text{break}, m_\text{1,low}, m_\text{high}) = \frac{1}{N}
	\begin{cases}
		  \left(\dfrac{m_1}{m_\text{break}}\right)^{-\alpha_1} & m_\text{1,low} \leq m_1  < m_\text{break} \\
		  \left(\dfrac{m_1}{m_\text{break}}\right)^{-\alpha_2} & m_\text{break} \leq m_1  < m_\text{high},
	\end{cases}
\end{eqnarray}

\noindent where $\alpha_1$ and $\alpha_2$ are the power law indices, the transition between the low-mass and high-mass power law occurs at $m_\text{break}$, and the normalization constant is

\begin{equation}
	N = m_{\rm break} \left[
    \frac{1 - \left(\tfrac{m_{1,\rm low}}{m_{\rm break}}\right)^{1-\alpha_1}}{1-\alpha_1}
    + \frac{\left(\tfrac{m_{1,\rm high}}{m_{\rm break}}\right)^{1-\alpha_2} - 1}{1-\alpha_2}\right].
\end{equation}

\noindent The full mixture distribution $\pi(m_1|{\PEhyperparam})$ is

\begin{align}
 \pi(m_1 | {\PEhyperparam})\propto & \Biggl[ \lambda_0 p_{\rm BP}(m_1|\alpha_1, \alpha_2, m_{\rm break}, m_{1,{\rm low}}, m_{\rm high}) + \lambda_1 \mathcal{N}_{lt}(m_1 | \mu_1, \sigma_1, \text{low} = m_{1,{\rm low}}) \\ \nonumber
& + (1 - \lambda_0 - \lambda_1) \mathcal{N}_{lt}(m_1|\mu_2, \sigma_2, \text{low} = m_{1,{\rm low}}) \Biggr] S(m_1 | m_{1,{\rm low}}, \delta_{m,1}),
\end{align}

\noindent where $\mathcal{N}_{lt}$ is a left-truncated normal distribution. The Planck tapering function $S$ ensures a smooth turn-on of the distribution in the range $(m_{\rm 1,low}, m_{\rm 1,low} + \delta_{m,1}]$ and is given by

\begin{eqnarray}
S(m|m_{\rm low}, \delta_m) =  \begin{cases}
0 & m < m_{\rm low}, \\\
[1 + f(m - m_{\rm low}, \delta_m)]^{-1}  & m_{\rm low} \leq m < m_{\rm low} + \delta_m, \\
1 & m_{\rm low} + \delta_m \leq m,
\end{cases}
\end{eqnarray}

\noindent with
$$
f(m', \delta_m) = \exp \left( \frac{\delta_m}{m'} + \frac{\delta_m}{m' - \delta_m} \right).
$$

We model the mass ratio as a power law with index $\beta_q$ and low-mass tapering applied to secondary mass $m_2$, conditioned on primary mass $m_1$,

\begin{equation}
p_{\rm PL}(q | m_1, \beta_q, m_{2,{\rm low}}, \delta_{m,2}) \propto q^{\beta_q} S(m_1 q | m_{2,{\rm low}}, \delta_{m,2}),
\end{equation}

\noindent with the same tapering function defined above. We numerically normalize the mass ratio distribution as a function of primary mass, and interpolate the normalization to arbitrary primary masses with a log-uniform grid across primary mass. We define $m_2 \leq m_1$, and therefore must enforce $m_{2, {\rm low}} \le m_{1, {\rm low}}$, which we enforce by the conditional priors described in Table \ref{tab:parameters_default}, alongisde the rest of the mass model priors.

\ExtendedBrokenPLTwoPeaks : This model incorporates correlations between the primary mass and mass ratio, by allowing each primary mass mixture component in the \BrokenPLTwoPeaks to be associated with a separate power law mass ratio model. In terms of the functions defined above, the model is expressed as

\begin{align} \label{eq:exbp2p}
\pi(m_1, q | {\PEhyperparam})\propto \Biggl[ &\lambda_0 p_{\rm BP}(m_1|\alpha_1, \alpha_2, m_{\rm break}, m_{1,{\rm low}}, m_{\rm high}) p_{\rm PL}(q | m_1, \beta_q^{\rm BP}, m_{2,{\rm low}}^{\rm BP}, \delta_{m,2}^{\rm BP})  \\ \nonumber
& + \lambda_1 \mathcal{N}_{lt}(m_1 | \mu_1, \sigma_1, \text{low} = m_{1,{\rm low}}) p_{\rm PL}(q | m_1, \beta_q^{\rm peak1}, m_{2,{\rm low}}^{\rm peak1}, \delta_{m,2}^{\rm peak1}) \\ \nonumber
&  + (1 - \lambda_0 - \lambda_1) \mathcal{N}_{lt}(m_1|\mu_2, \sigma_2, \text{low} = m_{1,{\rm low}}) p_{\rm PL}(q | m_1, \beta_q^{\rm peak2}, m_{2,{\rm low}}^{\rm peak2}, \delta_{m,2}^{\rm peak2}) \Biggr] \\ & \hspace{-4mm} \times S(m_1 | m_{1,{\rm low}}, \delta_{m,1}),\nonumber
\end{align}

\noindent where the superscripts (BP, peak1, and peak2) denote the different mass ratio hyperparameters ($\beta_q$, $m_{2,{\rm low}}$, $\delta_{m,2}$) for the broken power law, first Gaussian, and second Gaussian components. The priors on these hyperparameters are the same as those used in the \BrokenPLTwoPeaks model, listed in Table \ref{tab:parameters_default}, except for the power law index parameters $\beta_q$, which assume U$(-10,13)$.

\begin{deluxetable}{ccc}
\tablecaption{Summary of \BrokenPLTwoPeaks model parameters and priors.}
\tablehead{
\colhead{Parameter} & \colhead{Description} & \colhead{Prior}
}
\label{tab:parameters_default}
\startdata
$\alpha_1$ & Spectral index of 1st primary mass power law & U$(-4, 12)$ \\
$\alpha_2$ & Spectral index of 2nd primary mass power law & U$(-4, 12)$ \\
$m_{\rm break}$ & Power law break location & U$(20, 50)$ \\
$\mu_1$ & Location of the first peak & U$(5, 20)$ \\
$\sigma_1$ & Width of the first peak & U$(0,10)$ \\
$\mu_2$ & Location of the second peak & U$(25, 60)$ \\
$\sigma_2$ & Width of the second peak & U$(0,10)$ \\
$m_{\rm 1, low}$ & Lower edge of taper function & U$(3, 10)$ \\
$\delta_{\rm m,1}$ & Mass range of low mass tapering & U$(0,10)$ \\
$\lambda_0, \lambda_1$ & Mixing fractions between power law and peaks & ${\rm Dir}(\mathbf{\alpha}=(1,1,1))$ \\
$m_{\rm high}$ & Maximum mass for distribution, which is pinned & $\delta(m_{\rm high} - 300)$ \\
& to $m_{\rm high} = 300\,\Msun$ by default &  \\
\colrule
$\beta_q$ & Spectral index of mass ratio power law & U$(-2,7)$ \\
$m_{\rm 2, low}$ & Lower edge of taper function in $m_2$ & U$(3, m_{\rm 1, low})$ \\
$\delta_{\rm m,2}$ & Mass range of low mass tapering in $m_2$ & U$(0,10)$ \\
\enddata
\tablecomments{The priors for the \ExtendedBrokenPLTwoPeaks model are the same except for the $\beta_q$ parameters, which assume a prior of U$(-10,13)$.}
\end{deluxetable}

\subsection{\powerlawRedshift Model}\label{appendix:redshiftModels}

We model redshift evolution by the comoving merger rate density (Equation~\ref{eq:comoving_merger_rate}). In particular, we use the model
\begin{equation}
\pi(z|\kappa) \propto \frac{1}{1+z}\frac{\mathrm{d} V_c}{\mathrm{d} z}(1+z)^\kappa
\label{eq:powerlawredshift}
\end{equation}
where the prefactor converts from a rate density in comoving volume and source frame time to detector frame time and redshift. 
In other words, the comoving rate density scales as $\mathcal{R} \propto (1+z)^\kappa$.
We use a prior on $\kappa$ as specified in Table~\ref{tab:powerlawredshift}.

\begin{deluxetable}{ccc}
\tablecaption{Summary of \powerlawRedshift model parameter and prior.}
\tablehead{
\colhead{Parameter} & \colhead{Description} & \colhead{Prior}
}
\label{tab:powerlawredshift}
\startdata
$\kappa$ & Power-law index on comoving merger rate evolution & U($-10, 10$) \\
\enddata
\end{deluxetable}


%%%%%%%%%%%%%%%%%%%%%%%%%%%%%%%%%%%%%%%%%%%%%%%%%%%%%%%%
\subsection{Spin Models}\label{appendix:spinModels}
%%%%%%%%%%%%%%%%%%%%%%%%%%%%%%%%%%%%%%%%%%%%%%%%%%%%%%%%

BBH spins can be parameterized in several ways which are useful for GW data analysis. 
Here, we use dimensionless spin magnitudes for the \acp{BH} $\chi_1$ and $\chi_2$, and cosine tilt angles $\cos\theta_1$ and $\cos\theta_2$, as well as the effective spin parameters $\chieff$~\citep{Racine:2008qv,Ajith:2009bn,Damour:2001tu} and $\chip$~\citep{Schmidt:2010it,Schmidt:2012rh,Schmidt:2014iyl}. 
The effective inspiral spin $\chieff$ used because it is typically the most precisely measured \ac{BBH} spin parameter, due to its lower-order appearance post-Newtonian expansions~\citep{Arun:2008kb}. 
Other parametrizations of spin precession exist~\citep[e.g.,][]{Fairhurst:2019vut,Gerosa:2020aiw,Thomas:2020uqj} but for reasons of convention, we only work with the typical parameterization given in Equation 16 of \cite{GWTC:Introduction}.  
When modeling the effective spins, we use the analytic per-event parameter-estimation prior on $\chieff$, $\chip$, and $q$~\citep{Iwaya:2024zzq} in the calculation of the hierarchical likelihood (Equation~\ref{eq:hierarchical_likelihood}).
For all models presented in this work, we assume that azimuthal angles $\phi_i$ are distributed uniformly between $0$ and $2\pi$---the same as their parameter estimation prior---due to their typically uninformative individual-event posteriors.
See Table 3 of \cite{GWTC:Introduction} for more information about spin parameters.

We report results with spin vectors defined at the reference frequencies used for inference on each signal~\citep{GWTC:Methods}.
The effective spin $\chi_\mathrm{eff}$ is approximately conserved throughout the inspiral~\citep{Racine:2008qv,Gerosa:2015tea}, so its population distribution should be unaffected by the choice of reference frequency.
While the effective precessing spin $\chi_\mathrm{p}$ and the spin angles are dependent on reference frequency, they are comparatively weakly constrained by the data.
Additionally, their parameter estimation priors are invariant under time evolution, meaning measurements are robust against different reference frequencies.
At the current number of events and because of model dependence while measuring tilt distributions, we do not expect significant differences between the inferred tilt or $\chip$ distributions at different reference frequencies~\citep{Mould:2021xst}.
Our approach is consistent with past \ac{LVK} analyses, where evolved spins have not been used~\citep{LIGOScientific:2018jsj,LIGOScientific:2020kqk,KAGRA:2021duu}.


%%%%%%%%%%%%%%%%%%%%%%%%%%%%%%%%%%%%%%%%%%%%%%%%%%%%%%%%
\subsubsection{Component Spin Models}\label{appendix:componentSpinModels}
%%%%%%%%%%%%%%%%%%%%%%%%%%%%%%%%%%%%%%%%%%%%%%%%%%%%%%%%

\begin{deluxetable}{ccc}
\tablecaption{Summary of \default model parameters for spin magnitudes (Equation~\ref{eqn:spin_mag_model}) and tilt angles (Equation~\ref{eqn:spin_tilt_model},~\ref{eqn:spin_tilt_model_with_min}).}
\tablehead{
\colhead{Parameter} & \colhead{Description} & \colhead{Prior}
}
\label{tab:componentSpinPriors}
\startdata
$\mu_\chi$ & Location of the $\chi$ distribution & U($0$, $1$) \\
$\sigma_\chi$ & Width of the $\chi$ distribution & U($0.005$, $1$) \\
$\mu_t$ & Location of the Gaussian component of the $\cos\theta$ distribution & U($-1$, $1$) \\
$\sigma_t$ & Width of the Gaussian component of the $\cos\theta$ distribution & U($0.01$, $4$) \\
$\zeta$ & Fraction in the Gaussian component of the $\cos\theta$ distribution & U($0$, $1$) \\
\colrule
$t_{\rm min}$ & Minimum of the $\cos\theta$ distribution & U($-1$, $1$) \\
\enddata
\tablecomments{U stands for a uniform prior.}
\end{deluxetable}

\default:
We model the spin magnitudes ($\chi_i$) as a truncated Gaussian distribution between $0$ and $1$, assuming they are identically and independently distributed: 
%
\begin{equation}
    \pi(\chi_i|\mu_\chi,\sigma_\chi) = \mathcal{N}_{[0,1]}(\chi_1 | \mu_\chi, \sigma_\chi)\mathcal{N}_{[0,1]}(\chi_2 | \mu_\chi, \sigma_\chi)\,.
\label{eqn:spin_mag_model}
\end{equation}
%
This choice attempts to rectify shortcomings of the non-singular Beta distribution used to model the spin magnitudes in previous work~\citep{LIGOScientific:2018jsj, LIGOScientific:2020kqk, KAGRA:2021duu}; see Equation~\eqref{eqn:beta_distribution} below.
The non-singular Beta distribution is forced to go to $\pi(\chi)=0$ at $\chi=0,1$, which crucially does not allow for measurements of contributions to the population at $\chi=0$ or $\chi=1$, even though non-trivial contributions may exist~\citep{Callister:2022qwb,Galaudage:2021rkt,Hussain:2024qzl}.
Allowing the Beta distribution to be singular would add limited additional model flexibility, only adding the option for $\pi(\chi) = \infty$ at the $\chi=0,1$, but not anywhere in between $0$ and $\infty$. 
Thus, we opt for the truncated Gaussian model, which can take on a continuous range of values at $\chi$'s boundaries.

We model the distribution of the cosine spin tilt angle ($\cos\theta_i$) as a mixture between a Gaussian distribution truncated on $-1$ to $1$ and an isotropic distribution, assuming they are identically but \textit{not} independently distributed:
%
\begin{equation}
    \pi(\cos\theta_i| \mu_t, \sigma_t, \zeta ) = \zeta \,\mathcal{N}_{[-1,1]}(\cos\theta_1 | \mu_t, \sigma_t)\mathcal{N}_{[-1,1]}(\cos\theta_2 | \mu_t, \sigma_t) + \frac{1-\zeta}{4}\,.
\label{eqn:spin_tilt_model} 
\end{equation}
%
We here allow for the location of the Gaussian sub-population to vary, following~\cite{Vitale:2022dpa}, rather than fixing it at $\mu_t=1$ as was done in previous work~\citep{LIGOScientific:2018jsj, LIGOScientific:2020kqk, KAGRA:2021duu}.
Priors on the \default~hyperparameters are given in Table~\ref{tab:componentSpinPriors}.

{\sc Minimum Tilt Model}:
For the spin tilts, we also use a model in which $t_{\rm min}$ (where $t\equiv\cos\theta$), a hard lower truncation of $\pi(\cos\theta)$, is informed by the data. 
In this case the population distribution becomes 
\begin{multline}
    \pi(\cos\theta_i| \mu_t, \sigma_t, \zeta, t_{\rm min} ) \\ =
	\begin{cases} 
        \zeta \,\mathcal{N}_{[t_{\rm min},1]}(\cos\theta_1 | \mu_t, \sigma_t)\mathcal{N}_{[t_{\rm min},1]}(\cos\theta_2 | \mu_t, \sigma_t) + \displaystyle \frac{1-\zeta}{(1-t_{\rm min})^2} & \cos\theta_i > t_{\rm min}\,\,, \\
        0 & \cos\theta_i \leq t_{\rm min}\,\,.
	\end{cases}
	\label{eqn:spin_tilt_model_with_min}
\end{multline}
Equation \eqref{eqn:spin_tilt_model} is the same as Equation~\eqref{eqn:spin_tilt_model_with_min} with $t_{\rm min} = - 1$.
The prior on $t_{\rm min}$ and other hyperparameters are given in Table~\ref{tab:componentSpinPriors}.

{\sc Beta Distribution Spin Magnitude}:
In Appendix~\ref{appendix:model_comparison_spins}, we compare the \default~model to alternatives.
For the spin magnitudes, these are the Constrained and Unconstrained Beta Distributions, which both follow the form:
\begin{equation}
	\pi(\chi_i | \alpha, \beta) \propto {\chi_i}^{\alpha-1}{(1-\chi_i)}^{\beta-1} \,\,.
	\label{eqn:beta_distribution}
\end{equation}
The Constrained Beta was the {\sc Default} model for \gwtcthree~\citep{KAGRA:2021duu} and requires the shape parameters $\alpha, \beta > 1$, which forces $\pi(\chi_i=0,1) = 0$~\citep{Wysocki:2018mpo, LIGOScientific:2020kqk}.
The Unconstrained Beta relaxes this constraint on the the shape parameters. 
If $0< \alpha,\beta < 1$, the Beta distribution becomes singular, meaning $\pi(\chi_i=0,1) = \infty$.
Following \cite{KAGRA:2021duu}, we sample the mean and standard deviation of the Beta distribution, rather than $\alpha$ and $\beta$, and use the same priors for $\mu_\chi$ and $\sigma_\chi$ as given in Table~\ref{tab:componentSpinPriors} for the \default model.
For the Constrained Beta case, we additionally impose the cut $\alpha, \beta > 1$.
The relationship between $\{\mu_\chi,~\sigma_\chi\}$ and $\{\alpha,~\beta\}$ is given in Equation 5 of \cite{LIGOScientific:2018jsj}.


%%%%%%%%%%%%%%%%%%%%%%%%%%%%%%%%%%%%%%%%%%%%%%%%%%%%%%%%
\subsubsection{Identical versus Non-identically Distributed Spin Magnitudes and Tilts}
\label{appendix:iid}
%%%%%%%%%%%%%%%%%%%%%%%%%%%%%%%%%%%%%%%%%%%%%%%%%%%%%%%%

In Appendix~\ref{appendix:model_comparison_spins}, we investigate whether or not the primary and secondary spins are identically distributed, assuming the \default model.
\textit{Identical} distribution means that the two distributions share the same set of hyperparameters, while \textit{non-identical}  means that each is described by different hyperparameters.
Spin magnitudes are also \textit{independently} distributed, meaning that $\pi(\chi_1, \chi_2)$ is separable in terms of $\chi_1$ and $\chi_2$.
The acronym IID means they are independently and identically distributed, while IND means independently and non-identically distributed:
%
\begin{align}
	\chi_i~\mathrm{IID} & \implies \pi(\chi_{1,2} | \mu_\chi, \sigma_\chi) = \mathcal{N}_{[0,1]}(\chi_1 | \mu_\chi, \sigma_\chi) \, \mathcal{N}_{[0,1]}( \chi_2 | \mu_\chi, \sigma_\chi) \,\, , \label{eqn:iid_mags} \\
	\chi_i~\mathrm{IND} & \implies \pi(\chi_{1,2} | \mu_{\chi,1}, \sigma_{\chi,1}, \mu_{\chi,2}, \sigma_{\chi,2}) = \mathcal{N}_{[0,1]}(\chi_1 | \mu_{\chi,1}, \sigma_{\chi,1}) \, \mathcal{N}_{[0,1]}( \chi_2 | \mu_{\chi,2}, \sigma_{\chi,2}) \,\,.
	\label{eqn:ind_mags}
\end{align}
%
The tilt angles, on the other hand, are \textit{non-independently} distributed, meaning that the distribution $\pi(\cos\theta_1, \cos\theta_2)$ is not separable.
The acronym NID means they are non-independently but identically distributed, while NND means non-independently and non-identically distributed:
%
\begin{align}
	\cos\theta_i~\mathrm{NID} & \implies \pi(\cos\theta_{1,2}|\zeta, \mu_t, \sigma_t) = \zeta \,\mathcal{N}_{[-1,1]}(\cos\theta_1 | \mu_t, \sigma_t)\,\mathcal{N}_{[-1,1]}(\cos\theta_2 | \mu_t, \sigma_t) + \frac{1-\zeta}{4} \label{eqn:nid_tilts} \,\, ,\\
	\cos\theta_i~\mathrm{NND} & \implies \pi(\cos\theta_{1,2}|\zeta, \mu_{t,1}, \sigma_{t,1}, \mu_{t,2}, \sigma_{t,2}) \\ & \qquad\qquad\qquad\qquad = \zeta \,\mathcal{N}_{[-1,1]}(\cos\theta_1 | \mu_{t,1}, \sigma_{t,1})\,\mathcal{N}_{[-1,1]}(\cos\theta_2 | \mu_{t,2}, \sigma_{t,2})  + \frac{1-\zeta}{4} \,\,.
	\label{eqn:nnd_tilts}
\end{align}
%
For the non-identically distributed cases, each component's hyperparameters have the same priors as those listed in Table~\ref{tab:componentSpinPriors} for the identically distributed case. 
Table~\ref{tab:modelComparisonSpins} gives Bayes factors between different combinations of IID/IND spin magnitudes and NID/NND spin tilts; see associated dicsussion in Appendix~\ref{appendix:model_comparison_spins}.

%%%%%%%%%%%%%%%%%%%%%%%%%%%%%%%%%%%%%%%%%%%%%%%%%%%%%%%%
\subsubsection{Effective Spin Models}\label{appendix:effectiveSpinModels}
%%%%%%%%%%%%%%%%%%%%%%%%%%%%%%%%%%%%%%%%%%%%%%%%%%%%%%%

\begin{deluxetable}{cccc}
\tablecaption{Summary of \effectiveSpinModel~(Equation~\ref{eqn:chi_eff_chi_p_gaussian_model}) and \skewnormalChiEff~(Equation~\ref{eqn:chi_eff_skewnormal_model}) spin parameters.}
\tablehead{
\colhead{Parameter} & \colhead{Description} & \colhead{Prior G} & \colhead{Prior SN}
}
\label{tab:effectiveSpinPriors}
\startdata
$\mu_\mathrm{eff}$ & Location of the $\chieff$ distribution & U($-1$, $1$) & U($-1$, $1$) \\
$\sigma_\mathrm{eff}$ & Width of the $\chieff$ distribution & U($0.05$, $1$) & LU($0.01$, $4$) \\
$\mu_p$ & Location of the $\chi_\mathrm{p}$ distribution & U($0.05$, $1$) & U($0.01$, $1$) \\
$\sigma_p$ & Width of the $\chi_\mathrm{p}$ distribution & U($0.07$, $1$)& LU($0.01$, $1$)\\
\colrule
$\rho$ & Degree of correlation between $\chieff$ and $\chi_\mathrm{p}$ & U($-0.75$, $0.75$) & N/A \\
\colrule
$\epsilon $ & Skew of the $\chieff$ distribution & N/A & U($-1$, $1$) \\
\enddata
\tablecomments{The first column of priors ({\bf G}) gives those used for the \effectiveSpinModel results, the second column ({\bf SN}) is \skewnormalChiEff. U stands for a uniform prior; LU for log-uniform.}
\end{deluxetable}

\effectiveSpinModel: We here assume that the distribution of $\chieff$ and $\chip$ across the \ac{BBH} population is a bivariate truncated Gaussian~\citep{Miller:2020zox, Roulet:2018jbe} characterized by the location and width of the $\chieff$ and $\chip$ distributions, and the covariance between them:
\begin{equation}
    \pi(\chieff, \chip | \mathbf{\mu}, \mathbf{\Sigma}) \propto \mathcal{N}(\chieff, \chip | \mathbf{\mu}, \mathbf{\Sigma}) \,\, ,
\end{equation}
where $\mathbf \mu =(\mu_{\rm eff}, \mu_{\rm p})$ and
\begin{equation}
    \mathbf{\Sigma} =
    \begin{pmatrix}
        \sigma_{\rm eff}^2 & \rho\,\sigma_{\rm eff}\,\sigma_{\rm p} \\
        \rho\,\sigma_{\text{eff}}\,\sigma_{\rm p} & \sigma_{\rm p}^2
    \end{pmatrix}   \,\,.
\end{equation}
This expands to: 
\begin{equation}
    \pi(\chieff, \chip | \mathbf{\mu}, \mathbf{\Sigma}) \propto 
    \exp\left[ -\frac{1}{2(1 - \rho^2)}
    \left( \frac{(\chi_{\rm eff} - \mu_{\rm eff})^2}{\sigma_{\rm eff} ^2} 
    - \frac{2\rho (\chi_{\rm eff} - \mu_{\rm eff})(\chi_{\rm p} - \mu_{\rm p})}{\sigma_{\rm eff}\sigma_{\rm p}}
    + \frac{(\chi_{\rm p} - \mu_{\rm p})^2}{\sigma_{\rm p} ^2} \right) \right]\,\,.
    \label{eqn:chi_eff_chi_p_gaussian_model} 
\end{equation}
%
The bivariate Gaussian is truncated over the range $\chieff \in [-1,1]$, $\chip\in[0,1]$ and is normalized numerically.
Priors on the hyperparameters are given in Table \ref{tab:effectiveSpinPriors}.

\skewnormalChiEff: To account for observational evidence that the $\chieff$ distribution is not symmetric~\citep{Callister:2021fpo,Adamcewicz:2022hce,Banagiri:2025dxo}, we additionally model $\chieff$ as a skewed, truncated Gaussian distribution:

\begin{equation}
    \pi(\chieff|\mu_{\rm eff}, \sigma_{\rm eff}, \epsilon) \propto 
	\begin{cases} 
        (1 + \epsilon) \, \mathcal{N}_{[-1,1]}(\chieff | \mu_{\rm eff}, \sigma_{\rm eff} (1 + \epsilon)) \qquad \chieff \leq \mu_{\rm eff} \,\, ,\\ 
        (1 - \epsilon) \, \mathcal{N}_{[-1,1]}(\chieff | \mu_{\rm eff}, \sigma_{\rm eff} (1 - \epsilon)) \qquad \chieff \geq \mu_{\rm eff} \,\,.
	\end{cases}
    \label{eqn:chi_eff_skewnormal_model} 
\end{equation}
%
This model includes a parameter $\epsilon$ which describes the skew of the distribution, such that  $\epsilon > 0$ characterizes a distribution with more support for $\chieff<\mu_{\rm eff}$, while $\epsilon < 0$ has more support when $\chieff>\mu_{\rm eff}$. 
The \skewnormalChiEff~distribution reduces to a standard, symmetric Gaussian in the case that $\epsilon=0$.
We here model $\chip$ as a truncated normal distribution and infer its location and width. 
Priors on the hyperparameters are given in Table \ref{tab:effectiveSpinPriors}.

%%%%%%%%%%%%%%%%%%%%%%%%%%%%%%%%%%%%%%%%%%%%%%%%%%%%%%%%
\subsection{Copula Correlation Models}\label{appendix:copulaCorrelationModels}
%%%%%%%%%%%%%%%%%%%%%%%%%%%%%%%%%%%%%%%%%%%%%%%%%%%%%%%

\begin{deluxetable}{ccc}
\tablecaption{
    Summary of parameters exclusive to the \copulacorrelation, \linearcorrelation, and \splinecorrelation correlated models, along with their priors.
}
\tablehead{
    \colhead{Parameter} & \colhead{Description} & \colhead{Prior}
}
\label{tab:parameters_correlations}
\startdata
$\kappa_{x,y}$ & Level of correlation between parameters & U$(-20, 20)$ \\
& $x$ and $y$ inferred with a copula &
\\
\colrule
$\mu_{\mathrm{eff}|q}$ & $q=1$ intercept in linear $\mu_\mathrm{eff}(q)$ & U$(-1, 1)$ \\
$\delta \mu_{\mathrm{eff}|q}$ & Gradient in linear $\mu_\mathrm{eff}(q)$ & U$(-2, 2)$ \\
$\ln \sigma_{\mathrm{eff}|q}$ & $q=1$ intercept in linear $\ln \sigma_\mathrm{eff}(z)$ & U$(-5, 0)$ \\
$\delta \ln \sigma_{\mathrm{eff}|q}$ & Gradient in linear $\ln \sigma_\mathrm{eff}(q)$ & U$(-12, 4)$ \\
\colrule
$\mu_{\mathrm{eff}|z}$ & $z=0$ intercept in linear $\mu_\mathrm{eff}(z)$ & U$(-1, 1)$ \\
$\delta \mu_{\mathrm{eff}|z}$ & Gradient in linear $\mu_\mathrm{eff}(z)$ & U$(-1, 1)$ \\
$\ln \sigma_{\mathrm{eff}|z}$ & $z=0$ intercept in linear $\ln \sigma_\mathrm{eff}(z)$ & U$(-5, 0)$ \\
$\delta \ln \sigma_{\mathrm{eff}|z}$ & Gradient in linear $\ln \sigma_\mathrm{eff}(z)$ & U$(-3, 5)$ \\
\colrule
$\mu_{\mathrm{eff}|q}^i$ & $i^\mathrm{th}$ node in spline $\mu_\mathrm{eff}(q)$ & U$(-1, 1)$ \\
$\ln \sigma_{\mathrm{eff}|q}^i$ & $i^\mathrm{th}$ node in spline $\ln \sigma_\mathrm{eff}(q)$ & U$(-5, 0)$ \\
\colrule
$\mu_{\mathrm{eff}|z}^i$ & $i^\mathrm{th}$ node in spline $\mu_\mathrm{eff}(z)$ & U$(-1, 1)$ \\
$\ln \sigma_{\mathrm{eff}|z}^i$ & $i^\mathrm{th}$ node in spline $\ln \sigma_\mathrm{eff}(z)$ & U$(-5, 0)$ \\
\enddata
\tablecomments{    There are four variations of the \copulacorrelation~model in which $(x,y) = (q,\chieff)$, $(m_1,\chieff)$, $(z,\chieff)$ and finally, $(m_1,z)$.
    Within the spline models, we use four nodes $i$.
}
\end{deluxetable}

Copulas allow for a variable correlation between two parameters $x$ and $y$ while keeping the marginal distributions for the respective parameters fixed.
This is done using the fact that the cumulative distribution function (CDF) of any randomly distributed variable is itself a uniformly distributed variable between $0$ and $1$, in order to model the CDFs of $x$ and $y$ with a correlated two-dimensional uniform distribution.
This correlated two-dimensional uniform distribution, known as a copula density function, is dependent on a (hyper)parameter $\kappa_{x,y}$, which determines the level (or strength) of the correlation.
A coordinate transformation (which is determined by the chosen marginal distributions for $x$ and $y$) can then be applied to provide a correlated two-dimensional model in the desired coordinates.
Conveniently, the Jacobian for this transformation turns out to be the product of the chosen marginal distributions, meaning the two-dimensional population model can be written as
\begin{equation}
    \pi_{xy}(x,y|\PEhyperparam,\kappa_{x,y}) = \pi_c \left(u(x|\PEhyperparam),v(y|\PEhyperparam)|\kappa_{x,y}\right) \pi_x(x|\PEhyperparam) \pi_y(y|\PEhyperparam).
\end{equation}
Here, $\PEhyperparam$ are the set of hyperparameters governing the marginal distributions $\pi_x$ and $\pi_y$, while $\pi_c$ denotes the chosen copula density function.
Finally,
\begin{equation}
    u(x|\PEhyperparam) = \int_{x_{\min}}^x \mathrm{d} x' \ \pi_x(x'|\PEhyperparam),
\end{equation}
and
\begin{equation}
    v(y|\PEhyperparam) = \int_{y_{\min}}^y \mathrm{d} y' \ \pi_y(y'|\PEhyperparam),
\end{equation}
which we refer to as $u$ and $v$ for the sake of brevity, are the CDFs of $x$ and $y$ respectively.

The \copulacorrelation~models all assume the same marginal distributions and copula density functions, but differ in which pairs of parameters they correlate (i.e., which parameters $u$ and $v$ are functions of in the above notation).
Namely, the mass distribution is assumed to follow the default \BrokenPLTwoPeaks~model, redshift follows the default power-law redshift model, and $\chieff$ and $\chip$ are Gaussian distributed, but uncorrelated (i.e., they follow the \effectiveSpinModel~model with $\rho=0$).
We assume a Frank copula density function
\begin{equation}
    \pi_c\left(u, v | \kappa_{x,y}\right) = 
    \frac{-\kappa_{x,y} e^{-\kappa_{x,y} (u+v)} (e^{-\kappa_{x,y}} - 1)}{\left( e^{-\kappa_{x,y}} - e^{-\kappa_{x,y} u} - e^{-\kappa_{x,y} v} + e^{-\kappa_{x,y}(u+v)} \right)^2},
\end{equation}
where $\kappa_{x,y} \in (-\infty, \infty)$.
Positive values of $\kappa_{x,y}$ imply a correlation, while negative values of $\kappa_{x,y}$ imply an anti-correlation.
The Frank copula density function is not defined at $\kappa_{x,y} = 0$, but becomes uncorrelated as $\kappa_{x,y} \rightarrow 0$ from above and below, so we assume $\pi_c\left(u, v | \kappa_{x,y}=0\right) = 1$ (technically, making this a piece-wise function).
The Frank copula density function is chosen as it allows for positive and negative correlations that are symmetric about $\kappa_{x,y}=0$, and gives rise to correlated distributions that we believe appear physically reasonable \citep[][]{Adamcewicz:2022hce, Adamcewicz:2023mov}.
The prior on $\kappa_{x,y}$ (which is the same for all copula model variations), is given in Table~\ref{tab:parameters_correlations}.

Copulas are advantageous as they allow for a potential correlation to be quantified by a single variable $\kappa_{x,y}$, that otherwise has no influence on the distribution of the population.
However, relative to the linear and spline models explored below, they suffer from a lack of flexibility when it comes to covariance.
There are a limited number of copula density functions available, all of which introduce a correlation in a unique, but rigid way \citep[e.g,][]{Adamcewicz:2022hce, Adamcewicz:2023mov}.
Furthermore, known two-dimensional copulas depend on a single parameter $\kappa_{x,y}$, and cannot infer, for example, a separable broadening and correlation with the mean of a distribution simultaneously.
As a result, using copulas to probe for more complex structure in two-dimensions (assuming fixed marginal distributions), requires model comparison between a number of different copula density functions \citep[e.g.,][]{Adamcewicz:2023mov}.


%%%%%%%%%%%%%%%%%%%%%%%%%%%%%%%%%%%%%%%%%%%%%%%%%%%%%%%%
\subsection{Linear Correlation Models}\label{appendix:LinearCorrelationModels}
%%%%%%%%%%%%%%%%%%%%%%%%%%%%%%%%%%%%%%%%%%%%%%%%%%%%%%%
The $(q,\chieff)$ and $(z,\chieff)$ \linearcorrelation~models begin by assuming that $\chieff$ is Gaussian distributed for any given value of mass ratio $q$ and redshift $z$ respectively.
As the two models are otherwise identical, for the remainder of this Section, we substitute $x$ for $q$ and $z$.
This parameter-dependent Gaussian distribution for $\chieff$, $\pi(\chieff|x)$, is truncated at unphysical values of $|\chieff| \geq 1$.
From here, the mean and (natural log) width of the $\chieff$ distribution are allowed to evolve with the chosen variable $x$ linearly:
\begin{equation}
    \mu_\mathrm{eff}(x) = \mu_{\mathrm{eff}|x} + \delta \mu_{\mathrm{eff}|x} x,
\end{equation}
and
\begin{equation}
    \ln \sigma_\mathrm{eff}(x) = \ln \sigma_{\mathrm{eff}|x} + \delta \ln \sigma_{\mathrm{eff}|x_0} x.
\end{equation}
Here, $\mu_{\mathrm{eff}|x}$, $\delta \mu_{\mathrm{eff}|x}$, $\ln \sigma_{\mathrm{eff}|x}$, and $\delta \ln \sigma_{\mathrm{eff}|x_0}$ are all hyperparameters fit to the data, where $\delta \mu_{\mathrm{eff}|x}$ and $\delta \ln \sigma_{\mathrm{eff}|x_0}$ quantify the strength of the correlation between $x$ and the mean and width of the $\chieff$ distribution respectively.
The priors for these parameters are given in Table~\ref{tab:parameters_correlations}.
Meanwhile, masses follow the default \BrokenPLTwoPeaks~model and redshift is distributed according to the default power-law redshift model.

These models \citep[in the style of those presented in][]{Safarzadeh:2020mlb, Callister:2021fpo, Biscoveanu:2022qac}, therefore allow us to infer separable trends in the mean and width of the $\chieff$ distribution with another parameter.
The inferred correlations should be interpreted with care as, unlike the copula models defined in Section~\ref{appendix:copulaCorrelationModels}, the correlation hyperparameters here will also affect the shape of the marginal $\chieff$ distribution.
As such, it can be difficult to be certain whether an inferred correlation is entirely due to a trend between $\chieff$ and another parameter, or is in part due to a better fit to the marginal $\chieff$ distribution \citep{Adamcewicz:2022hce}.


%%%%%%%%%%%%%%%%%%%%%%%%%%%%%%%%%%%%%%%%%%%%%%%%%%%%%%%%
\subsection{Spline Correlation Models}\label{appendix:SplineCorrelationModels}
%%%%%%%%%%%%%%%%%%%%%%%%%%%%%%%%%%%%%%%%%%%%%%%%%%%%%%%
The \splinecorrelation~models~\citep{Heinzel:2023hlb} are similar to the \linearcorrelation~models defined in Section~\ref{appendix:LinearCorrelationModels}.
The key difference is that the \splinecorrelation~models more flexibly model the parameter-dependent mean $\mu_\mathrm{eff}(x)$ and (natural log) width $\ln \sigma_\mathrm{eff}(x)$ of the $\chieff$ distribution as cubic splines dependent on $x$.
Each spline has four nodes $\mu_\mathrm{eff}^i$ and $\ln \sigma_{\mathrm{eff}|x}^i$ that are placed uniformly in $\log_{10} x$, and are inferred from the data.
The priors for these nodes are given in Table~\ref{tab:parameters_correlations}.
